\setcounter{chapter}{1}
\pagestyle{fancy}

\chapter{行列式}

本章前六节完全对应原讲义中的专题一与教材相关部分，原讲义以习题训练为主，因此这里保留其题目顺序，并在答案版中补充相应解答。若题目未另作说明，本章默认要求计算所给行列式。


\section{大对角形及递推法}

\begin{example}
计算
\[
D_n=
\begin{vmatrix}
a & b & 0 & \cdots & 0\\
c & a & b & \ddots & \vdots\\
0 & c & a & \ddots & 0\\
\vdots & \ddots & \ddots & \ddots & b\\
0 & \cdots & 0 & c & a
\end{vmatrix}.
\]
\end{example}
\exampleblank

\begin{solution}
按最后一行或最后一列展开可得
\[
D_n=aD_{n-1}-bcD_{n-2},
\qquad
D_0=1,\quad D_1=a.
\]
特征方程为
\[
\lambda^2-a\lambda+bc=0.
\]
记
\[
\lambda_{\pm}
=
\dfrac{a\pm\sqrt{a^2-4bc}}{2}.
\]
若 $a^2\neq4bc$，则
\[
D_n
=
\dfrac{\lambda_+^{\,n+1}-\lambda_-^{\,n+1}}
{\lambda_+-\lambda_-}.
\]
若 $a^2=4bc$，令 $\lambda=\dfrac{a}{2}$，则
\[
D_n=(n+1)\lambda^n.
\]
\end{solution}


\begin{example}
计算
\[
D_n=
\begin{vmatrix}
\alpha+\beta & \alpha\beta & 0 & \cdots & 0\\
1 & \alpha+\beta & \alpha\beta & \ddots & \vdots\\
0 & 1 & \alpha+\beta & \ddots & 0\\
\vdots & \ddots & \ddots & \ddots & \alpha\beta\\
0 & \cdots & 0 & 1 & \alpha+\beta
\end{vmatrix},
\qquad
\alpha\neq\beta.
\]
\end{example}
\exampleblank

\begin{solution}
递推关系为
\[
D_n=(\alpha+\beta)D_{n-1}-\alpha\beta D_{n-2},
\qquad
D_0=1,\quad D_1=\alpha+\beta.
\]
特征方程
\[
\lambda^2-(\alpha+\beta)\lambda+\alpha\beta=0
\]
的两根为 $\alpha,\beta$，故
\[
D_n
=
\dfrac{\alpha^{n+1}-\beta^{n+1}}{\alpha-\beta}.
\]
\end{solution}


\begin{example}
计算
\[
D_n=
\begin{vmatrix}
\cos x & 1 & 0 & \cdots & 0\\
1 & 2\cos x & 1 & \ddots & \vdots\\
0 & 1 & 2\cos x & \ddots & 0\\
\vdots & \ddots & \ddots & \ddots & 1\\
0 & \cdots & 0 & 1 & 2\cos x
\end{vmatrix}.
\]
\end{example}
\exampleblank

\begin{solution}
有
\[
D_1=\cos x,\qquad
D_2=2\cos^2x-1=\cos2x,
\]
且
\[
D_n=2\cos x\,D_{n-1}-D_{n-2}.
\]
由
\[
\cos nx
=
2\cos x\cos((n-1)x)-\cos((n-2)x)
\]
可知
\[
\boxed{D_n=\cos nx}.
\]
\end{solution}


\begin{example}
计算
\[
D_n=
\begin{vmatrix}
2n & n & 0 & \cdots & 0\\
n & 2n & n & \ddots & \vdots\\
0 & n & 2n & \ddots & 0\\
\vdots & \ddots & \ddots & \ddots & n\\
0 & \cdots & 0 & n & 2n
\end{vmatrix}.
\]
\end{example}
\exampleblank

\begin{solution}
从每一行提出 $n$，得
\[
D_n
=
n^n
\begin{vmatrix}
2 & 1 & 0 & \cdots & 0\\
1 & 2 & 1 & \ddots & \vdots\\
0 & 1 & 2 & \ddots & 0\\
\vdots & \ddots & \ddots & \ddots & 1\\
0 & \cdots & 0 & 1 & 2
\end{vmatrix}.
\]
设后一个行列式为 $E_n$，则
\[
E_n=2E_{n-1}-E_{n-2},
\qquad
E_1=2,\quad E_2=3,
\]
故 $E_n=n+1$。因此
\[
\boxed{D_n=(n+1)n^n}.
\]
\end{solution}


\begin{example}
计算
\[
D_n=
\begin{vmatrix}
2 & 1 & 0 & \cdots & 0\\
1 & 2 & 1 & \ddots & \vdots\\
0 & 1 & 2 & \ddots & 0\\
\vdots & \ddots & \ddots & \ddots & 1\\
0 & \cdots & 0 & 1 & 2
\end{vmatrix}.
\]
\end{example}
\exampleblank

\begin{solution}
递推
\[
D_n=2D_{n-1}-D_{n-2},
\qquad
D_1=2,\quad D_2=3.
\]
由归纳法立即得到
\[
\boxed{D_n=n+1}.
\]
\end{solution}


\begin{example}
计算
\[
D_n=
\begin{vmatrix}
1-a_1 & a_2 & 0 & \cdots & 0\\
-1 & 1-a_2 & a_3 & \ddots & \vdots\\
0 & -1 & 1-a_3 & \ddots & 0\\
\vdots & \ddots & \ddots & \ddots & a_n\\
0 & \cdots & 0 & -1 & 1-a_n
\end{vmatrix}.
\]
\end{example}
\exampleblank

\begin{solution}
记 $D_0=1$，则
\[
D_1=1-a_1,
\]
并有
\[
D_k=(1-a_k)D_{k-1}+a_kD_{k-2}.
\]
因此
\[
D_k-D_{k-1}
=
-a_k(D_{k-1}-D_{k-2}).
\]
由
\[
D_1-D_0=-a_1
\]
递推得
\[
D_k-D_{k-1}
=
(-1)^k a_1a_2\cdots a_k.
\]
累加即得
\[
\boxed{
D_n
=
1-a_1+a_1a_2-a_1a_2a_3+\cdots
+(-1)^n a_1a_2\cdots a_n
}.
\]
\end{solution}


\begin{example}
设
\[
|A_n|=
\begin{vmatrix}
a_0+a_1 & a_1 & 0 & \cdots & 0\\
a_1 & a_1+a_2 & a_2 & \ddots & \vdots\\
0 & a_2 & a_2+a_3 & \ddots & 0\\
\vdots & \ddots & \ddots & \ddots & a_{n-1}\\
0 & \cdots & 0 & a_{n-1} & a_{n-1}+a_n
\end{vmatrix}.
\]
证明：
\[
|A_n|
=
a_0a_1\cdots a_n
\left(
\dfrac{1}{a_0}
+\dfrac{1}{a_1}
+\cdots
+\dfrac{1}{a_n}
\right).
\]
\end{example}
\exampleblank

\begin{solution}
记右端为
\[
F_n=\sum\limits_{k=0}^{n}\prod_{\substack{0\leq j\leq n\\j\neq k}}a_j.
\]
按最后一行展开三对角行列式可得
\[
|A_n|
=
(a_{n-1}+a_n)|A_{n-1}|
-a_{n-1}^2|A_{n-2}|.
\]
直接代入 $F_{n-1},F_{n-2}$ 可验证
\[
F_n
=
(a_{n-1}+a_n)F_{n-1}
-a_{n-1}^2F_{n-2},
\]
而 $n=1$ 时结论显然成立，故由归纳法得
\[
\boxed{
|A_n|
=
a_0a_1\cdots a_n
\sum\limits_{k=0}^{n}\dfrac{1}{a_k}
}.
\]
若某个 $a_k=0$，将右端理解为前面的无分母多项式形式即可。
\end{solution}


\begin{example}
已知 $a_i>0,b_i>0\ (i=1,2,\ldots,n)$，证明
\[
D_n=
\begin{vmatrix}
a_1 & b_1 & 0 & \cdots & 0\\
-1 & a_2 & b_2 & \ddots & \vdots\\
0 & -1 & a_3 & \ddots & 0\\
\vdots & \ddots & \ddots & \ddots & b_{n-1}\\
0 & \cdots & 0 & -1 & a_n
\end{vmatrix}>0.
\]
\end{example}
\exampleblank

\begin{solution}
递推关系为
\[
D_n=a_nD_{n-1}+b_{n-1}D_{n-2},
\qquad
D_0=1,\quad D_1=a_1>0.
\]
由于 $a_n>0,b_{n-1}>0$，由归纳法有
\[
D_n>0
\]
对一切正整数 $n$ 成立。
\end{solution}


\newpage

\section{行（列）和相等的行列式}

\begin{example}
计算
\[
D_n=
\begin{vmatrix}
0 & 1 & \cdots & 1\\
1 & 0 & \cdots & 1\\
\vdots & \vdots & \ddots & \vdots\\
1 & 1 & \cdots & 0
\end{vmatrix}.
\]
\end{example}
\exampleblank

\begin{solution}
矩阵可写成
\[
J-I_n,
\]
其中 $J$ 为全 $1$ 矩阵。向量 $(1,\ldots,1)^{\mathsf T}$ 对应特征值
$n-1$，而其正交补上的特征值均为 $-1$。故
\[
\boxed{
D_n=(-1)^{n-1}(n-1)
}.
\]
\end{solution}


\begin{example}
计算
\[
D_n=
\begin{vmatrix}
a_1+b & a_2 & a_3 & \cdots & a_n\\
a_1 & a_2+b & a_3 & \cdots & a_n\\
a_1 & a_2 & a_3+b & \cdots & a_n\\
\vdots & \vdots & \vdots & \ddots & \vdots\\
a_1 & a_2 & a_3 & \cdots & a_n+b
\end{vmatrix}.
\]
\end{example}
\exampleblank

\begin{solution}
记
\[
\boldsymbol{1}=(1,\ldots,1)^{\mathsf T},
\qquad
a=(a_1,\ldots,a_n)^{\mathsf T}.
\]
则
\[
A=bI_n+\boldsymbol{1}a^{\mathsf T}.
\]
由秩一行列式公式
\[
\det(I+uv^{\mathsf T})=1+v^{\mathsf T}u
\]
可得
\[
\boxed{
D_n=b^{\,n-1}
\left(b+\sum\limits_{i=1}^{n}a_i\right)
}.
\]
该式作为关于 $b$ 的多项式恒等式，对 $b=0$ 也成立。
\end{solution}


\begin{example}
计算
\[
\begin{vmatrix}
5&6&7&8\\
6&7&8&5\\
7&8&5&6\\
8&5&6&7
\end{vmatrix}.
\]
\end{example}
\exampleblank

\begin{solution}
这是四阶循环行列式。取四次单位根
\[
1,\quad i,\quad -1,\quad -i,
\]
相应特征值为
\[
26,\quad -2-2i,\quad -2,\quad -2+2i.
\]
故
\[
D
=
26(-2)(-2-2i)(-2+2i)
=
\boxed{-416}.
\]
\end{solution}


\begin{example}
计算
\[
D_n=
\begin{vmatrix}
1&1&\cdots&1&2-n\\
1&1&\cdots&2-n&1\\
\vdots&\vdots&\reflectbox{$\ddots$}&\vdots&\vdots\\
1&2-n&\cdots&1&1\\
2-n&1&\cdots&1&1
\end{vmatrix}.
\]
\end{example}
\exampleblank

\begin{solution}
设 $P$ 为反序置换矩阵，则
\[
A=J-(n-1)P.
\]
在向量 $\boldsymbol{1}$ 上，
\[
A\boldsymbol{1}=\boldsymbol{1},
\]
故该方向的特征值为 $1$。在
\[
\boldsymbol{1}^{\perp}
\]
上有 $J=0$，于是 $A=-(n-1)P$。又
\[
\det P=(-1)^{\frac{n(n-1)}{2}}.
\]
因此
\[
\boxed{
D_n
=
(-1)^{n-1+\frac{n(n-1)}{2}}(n-1)^{n-1}
=
(-1)^{\frac{(n-1)(n+2)}{2}}(n-1)^{n-1}
}.
\]
\end{solution}


\newpage

\section{升阶法}

为方便书写，记
\[
\Delta(x_1,\ldots,x_n)
=
\prod\limits_{1\leq i<j\leq n}(x_j-x_i),
\]
并以 $e_k=e_k(x_1,\ldots,x_n)$ 表示第 $k$ 个初等对称多项式，约定 $e_0=1$。

\begin{example}
计算
\[
D_n=
\begin{vmatrix}
1+x_1&1+x_1^2&\cdots&1+x_1^n\\
1+x_2&1+x_2^2&\cdots&1+x_2^n\\
\vdots&\vdots&&\vdots\\
1+x_n&1+x_n^2&\cdots&1+x_n^n
\end{vmatrix}.
\]
\end{example}
\exampleblank

\begin{solution}
将各列写成“常数列 $1$”与幂列之和，并按列作多线性展开。由于两列同时取常数列时行列式为零，只保留“没有常数列”和“恰有一列取常数列”的项。利用缺一阶广义 Vandermonde 行列式，可得
\[
\boxed{
D_n
=
\Delta(x_1,\ldots,x_n)
\left(
e_n+e_{n-1}-e_{n-2}+e_{n-3}-\cdots+(-1)^{n-1}
\right)
}.
\]
\end{solution}


\begin{example}
计算
\[
D_n=
\begin{vmatrix}
1+x&x&\cdots&x\\
x&2+x&\cdots&x\\
\vdots&\vdots&\ddots&\vdots\\
x&x&\cdots&n+x
\end{vmatrix}.
\]
\end{example}
\exampleblank

\begin{solution}
写成
\[
A=\diag(1,2,\ldots,n)+xJ.
\]
由矩阵行列式引理，
\[
\det(D+x\boldsymbol{1}\boldsymbol{1}^{\mathsf T})
=
\det D
\left(
1+x\boldsymbol{1}^{\mathsf T}D^{-1}\boldsymbol{1}
\right),
\]
故
\[
\boxed{
D_n
=
n!
\left(
1+x\sum\limits_{k=1}^{n}\dfrac{1}{k}
\right)
}.
\]
\end{solution}


\begin{example}
计算
\[
D_n=
\begin{vmatrix}
a_1+b_1&a_2&a_3&\cdots&a_n\\
a_1&a_2+b_2&a_3&\cdots&a_n\\
a_1&a_2&a_3+b_3&\cdots&a_n\\
\vdots&\vdots&\vdots&\ddots&\vdots\\
a_1&a_2&a_3&\cdots&a_n+b_n
\end{vmatrix}.
\]
\end{example}
\exampleblank

\begin{solution}
有
\[
A=\diag(b_1,\ldots,b_n)+\boldsymbol{1}a^{\mathsf T}.
\]
于是
\[
\boxed{
D_n
=
\prod\limits_{i=1}^{n}b_i
+
\sum\limits_{i=1}^{n}
a_i
\prod_{\substack{1\leq j\leq n\\j\neq i}}b_j
}.
\]
当各 $b_i\neq0$ 时，也可写成
\[
D_n
=
\left(\prod\limits_{i=1}^{n}b_i\right)
\left(
1+\sum\limits_{i=1}^{n}\dfrac{a_i}{b_i}
\right).
\]
\end{solution}


\begin{example}
计算
\[
D_n=
\begin{vmatrix}
1&x_1&x_1^2&\cdots&x_1^{n-2}&x_1^n\\
1&x_2&x_2^2&\cdots&x_2^{n-2}&x_2^n\\
\vdots&\vdots&\vdots&&\vdots&\vdots\\
1&x_n&x_n^2&\cdots&x_n^{n-2}&x_n^n
\end{vmatrix}.
\]
\end{example}
\exampleblank

\begin{solution}
这是 Vandermonde 行列式中缺少 $x^{n-1}$ 列的情形。由广义 Vandermonde 公式，
\[
\boxed{
D_n
=
(x_1+x_2+\cdots+x_n)
\Delta(x_1,\ldots,x_n)
}.
\]
\end{solution}


\begin{example}
计算
\[
D_n=
\begin{vmatrix}
1&x_1(x_1-1)&x_1^2(x_1-1)&\cdots&x_1^{n-1}(x_1-1)\\
1&x_2(x_2-1)&x_2^2(x_2-1)&\cdots&x_2^{n-1}(x_2-1)\\
\vdots&\vdots&\vdots&&\vdots\\
1&x_n(x_n-1)&x_n^2(x_n-1)&\cdots&x_n^{n-1}(x_n-1)
\end{vmatrix}.
\]
\end{example}
\exampleblank

\begin{solution}
第 $k+1$ 列对应多项式
\[
x^k(x-1)=x^{k+1}-x^k,
\qquad k=1,\ldots,n-1.
\]
逐列消去低次项后，可化为缺少一次项的广义 Vandermonde 型。结果为
\[
\boxed{
D_n
=
\Delta(x_1,\ldots,x_n)
\left(
e_{n-1}-e_{n-2}+e_{n-3}-\cdots+(-1)^{n-1}
\right)
}.
\]
\end{solution}


\begin{example}
计算
\[
D_n=
\begin{vmatrix}
1&x_1(x_1-a)&x_1^2(x_1-a)&\cdots&x_1^{n-1}(x_1-a)\\
1&x_2(x_2-a)&x_2^2(x_2-a)&\cdots&x_2^{n-1}(x_2-a)\\
\vdots&\vdots&\vdots&&\vdots\\
1&x_n(x_n-a)&x_n^2(x_n-a)&\cdots&x_n^{n-1}(x_n-a)
\end{vmatrix}.
\]
\end{example}
\exampleblank

\begin{solution}
与上一题同理，利用
\[
x^k(x-a)=x^{k+1}-ax^k
\]
逐列化简，得到
\[
\boxed{
D_n
=
\Delta(x_1,\ldots,x_n)
\left(
e_{n-1}
-ae_{n-2}
+a^2e_{n-3}
-\cdots
+(-a)^{n-1}
\right)
}.
\]
\end{solution}


\newpage

\section{\texorpdfstring{$A(t)$}{A(t)} 的应用}

\begin{example}
计算
\[
D_n=
\begin{vmatrix}
x&b&b&\cdots&b\\
a&x&b&\cdots&b\\
a&a&x&\cdots&b\\
\vdots&\vdots&\vdots&\ddots&\vdots\\
a&a&a&\cdots&x
\end{vmatrix}.
\]
\end{example}
\exampleblank

\begin{solution}
当 $a\neq b$ 时，连续作相邻行差或使用参数行列式可得
\[
\boxed{
D_n
=
\dfrac{
a(x-b)^n-b(x-a)^n
}{a-b}
}.
\]
当 $a=b$ 时取极限，得
\[
\boxed{
D_n
=
(x-a)^{n-1}\bigl(x+(n-1)a\bigr)
}.
\]
\end{solution}


\begin{example}
计算
\[
D_n=
\begin{vmatrix}
x_1&y&y&\cdots&y\\
z&x_2&y&\cdots&y\\
z&z&x_3&\cdots&y\\
\vdots&\vdots&\vdots&\ddots&\vdots\\
z&z&z&\cdots&x_n
\end{vmatrix},
\qquad
y\neq z.
\]
\end{example}
\exampleblank

\begin{solution}
利用相邻行差，或对上一题作变参数推广，可得
\[
\boxed{
D_n
=
\dfrac{
y\prod\limits_{i=1}^{n}(x_i-z)
-
z\prod\limits_{i=1}^{n}(x_i-y)
}{y-z}
}.
\]
\end{solution}


\begin{example}
设 $A=(a_{ij})$，$A_{ij}$ 表示元素 $a_{ij}$ 的代数余子式。证明
\[
\begin{vmatrix}
a_{11}-a_{12}&a_{12}-a_{13}&\cdots&a_{1,n-1}-a_{1n}&1\\
a_{21}-a_{22}&a_{22}-a_{23}&\cdots&a_{2,n-1}-a_{2n}&1\\
\vdots&\vdots&&\vdots&\vdots\\
a_{n1}-a_{n2}&a_{n2}-a_{n3}&\cdots&a_{n,n-1}-a_{nn}&1
\end{vmatrix}
=
\sum\limits_{i,j=1}^{n}A_{ij}.
\]
\end{example}
\exampleblank

\begin{solution}
记 $C_k$ 为 $A$ 的第 $k$ 列。左端矩阵的前 $n-1$ 列为
\[
C_1-C_2,\ C_2-C_3,\ \ldots,\ C_{n-1}-C_n,
\]
最后一列为 $\boldsymbol{1}$。

先设 $A$ 可逆。于是左端等于
\[
\det A\,
\det
\begin{pmatrix}
e_1-e_2& e_2-e_3&\cdots&e_{n-1}-e_n&A^{-1}\boldsymbol{1}
\end{pmatrix}.
\]
对任意列向量 $v$，
\[
\det
\begin{pmatrix}
e_1-e_2&\cdots&e_{n-1}-e_n&v
\end{pmatrix}
=
\sum\limits_{i=1}^{n}v_i.
\]
故原式等于
\[
\det A\,
\boldsymbol{1}^{\mathsf T}A^{-1}\boldsymbol{1}
=
\boldsymbol{1}^{\mathsf T}\operatorname{adj}(A)\boldsymbol{1}
=
\sum\limits_{i,j=1}^{n}A_{ij}.
\]
两端均为 $a_{ij}$ 的多项式，故由多项式恒等原理，对奇异 $A$ 也成立。
\end{solution}


\begin{example}
计算
\[
D_n=
\begin{vmatrix}
1+x_1&1+x_1^2&\cdots&1+x_1^n\\
1+x_2&1+x_2^2&\cdots&1+x_2^n\\
\vdots&\vdots&&\vdots\\
1+x_n&1+x_n^2&\cdots&1+x_n^n
\end{vmatrix}.
\]
\end{example}
\exampleblank

\begin{solution}
这与例题 2.3.1 相同，因此
\[
\boxed{
D_n
=
\Delta(x_1,\ldots,x_n)
\left(
e_n+e_{n-1}-e_{n-2}+e_{n-3}-\cdots+(-1)^{n-1}
\right)
}.
\]
\end{solution}


\newpage

\section{（大小）拆分法}

\subsection{大拆分法}

\begin{example}[\markstar]
设 $A=(a_{ij})$，已知 $|A|$ 及各代数余子式 $A_{ij}$，计算
\[
|A(t)|
=
\begin{vmatrix}
a_{11}+t&a_{12}+t&\cdots&a_{1n}+t\\
a_{21}+t&a_{22}+t&\cdots&a_{2n}+t\\
\vdots&\vdots&&\vdots\\
a_{n1}+t&a_{n2}+t&\cdots&a_{nn}+t
\end{vmatrix}.
\]
\end{example}
\exampleblank

\begin{solution}
有
\[
A(t)=A+t\boldsymbol{1}\boldsymbol{1}^{\mathsf T}.
\]
由行列式关于任一列的多线性知，展开时至多有一列取 $t\boldsymbol{1}$，故 $|A(t)|$ 关于 $t$ 至多一次。常数项为 $|A|$，一次项为全体代数余子式之和，因此
\[
\boxed{
|A(t)|
=
|A|
+t\sum\limits_{i,j=1}^{n}A_{ij}
}.
\]
\end{solution}


\begin{example}
计算
\[
D_n=
\begin{vmatrix}
a_1+b_1&a_1+b_2&\cdots&a_1+b_n\\
a_2+b_1&a_2+b_2&\cdots&a_2+b_n\\
\vdots&\vdots&&\vdots\\
a_n+b_1&a_n+b_2&\cdots&a_n+b_n
\end{vmatrix}.
\]
\end{example}
\exampleblank

\begin{solution}
矩阵可写成
\[
a\boldsymbol{1}^{\mathsf T}
+
\boldsymbol{1}b^{\mathsf T},
\]
故秩不超过 $2$。因此当 $n\geq3$ 时
\[
\boxed{D_n=0}.
\]
另外，
\[
D_1=a_1+b_1,
\]
而 $n=2$ 时
\[
D_2=(a_1-a_2)(b_2-b_1).
\]
\end{solution}


\begin{example}
计算
\[
D_n=
\begin{vmatrix}
a&b&\cdots&b\\
b&a&\cdots&b\\
\vdots&\vdots&\ddots&\vdots\\
b&b&\cdots&a
\end{vmatrix}.
\]
\end{example}
\exampleblank

\begin{solution}
写成
\[
A=(a-b)I_n+bJ.
\]
故特征值为
\[
a+(n-1)b
\]
一次，以及
\[
a-b
\]
共 $n-1$ 次。因此
\[
\boxed{
D_n=(a-b)^{n-1}\bigl(a+(n-1)b\bigr)
}.
\]
\end{solution}


\begin{example}
计算
\[
D_n=
\begin{vmatrix}
x_1&a_2&a_3&\cdots&a_n\\
a_1&x_2&a_3&\cdots&a_n\\
a_1&a_2&x_3&\cdots&a_n\\
\vdots&\vdots&\vdots&\ddots&\vdots\\
a_1&a_2&a_3&\cdots&x_n
\end{vmatrix},
\qquad
a_i\neq x_i.
\]
\end{example}
\exampleblank

\begin{solution}
写成
\[
A
=
\diag(x_1-a_1,\ldots,x_n-a_n)
+
\boldsymbol{1}a^{\mathsf T}.
\]
由矩阵行列式引理，
\[
\boxed{
D_n
=
\left(
\prod\limits_{i=1}^{n}(x_i-a_i)
\right)
\left(
1+
\sum\limits_{i=1}^{n}\dfrac{a_i}{x_i-a_i}
\right)
}.
\]
\end{solution}


\begin{example}
计算
\[
D_n=
\begin{vmatrix}
1+x&x&\cdots&x\\
x&2+x&\cdots&x\\
\vdots&\vdots&\ddots&\vdots\\
x&x&\cdots&n+x
\end{vmatrix}.
\]
\end{example}
\exampleblank

\begin{solution}
与例题 2.3.2 相同，
\[
\boxed{
D_n
=
n!
\left(
1+x\sum\limits_{k=1}^{n}\dfrac{1}{k}
\right)
}.
\]
\end{solution}


\begin{example}
计算
\[
D_n=
\begin{vmatrix}
a_1+b_1&a_2&a_3&\cdots&a_n\\
a_1&a_2+b_2&a_3&\cdots&a_n\\
a_1&a_2&a_3+b_3&\cdots&a_n\\
\vdots&\vdots&\vdots&\ddots&\vdots\\
a_1&a_2&a_3&\cdots&a_n+b_n
\end{vmatrix}.
\]
\end{example}
\exampleblank

\begin{solution}
与例题 2.3.3 相同，
\[
\boxed{
D_n
=
\prod\limits_{i=1}^{n}b_i
+
\sum\limits_{i=1}^{n}
a_i
\prod_{\substack{1\leq j\leq n\\j\neq i}}b_j
}.
\]
\end{solution}


\subsection{小拆分法}

\begin{example}
计算
\[
D_n=
\begin{vmatrix}
x&a&a&\cdots&a\\
b&x&a&\cdots&a\\
b&b&x&\cdots&a\\
\vdots&\vdots&\vdots&\ddots&\vdots\\
b&b&b&\cdots&x
\end{vmatrix}.
\]
\end{example}
\exampleblank

\begin{solution}
当 $a\neq b$ 时，由例题 2.4.1，
\[
\boxed{
D_n
=
\dfrac{
a(x-b)^n-b(x-a)^n
}{a-b}
}.
\]
当 $a=b$ 时，
\[
\boxed{
D_n=(x-a)^{n-1}\bigl(x+(n-1)a\bigr)
}.
\]
\end{solution}


\begin{example}
计算
\[
D_n=
\begin{vmatrix}
x_1&a&a&\cdots&a\\
b&x_2&a&\cdots&a\\
b&b&x_3&\cdots&a\\
\vdots&\vdots&\vdots&\ddots&\vdots\\
b&b&b&\cdots&x_n
\end{vmatrix}.
\]
\end{example}
\exampleblank

\begin{solution}
若 $a\neq b$，则由例题 2.4.2，
\[
\boxed{
D_n
=
\dfrac{
a\prod\limits_{i=1}^{n}(x_i-b)
-
b\prod\limits_{i=1}^{n}(x_i-a)
}{a-b}
}.
\]
当 $a=b$ 时可由上式取极限得到相应结果。
\end{solution}


\newpage

\section{矩阵分解技巧与降阶公式（打洞原理）}

\begin{example}
计算
\[
D_n=
\begin{vmatrix}
a_1b_1&a_1b_2&\cdots&a_1b_n\\
a_2b_1&a_2b_2&\cdots&a_2b_n\\
\vdots&\vdots&&\vdots\\
a_nb_1&a_nb_2&\cdots&a_nb_n
\end{vmatrix}.
\]
\end{example}
\exampleblank

\begin{solution}
矩阵为外积
\[
ab^{\mathsf T},
\]
其秩不超过 $1$。故当 $n\geq2$ 时
\[
\boxed{D_n=0}.
\]
$n=1$ 时 $D_1=a_1b_1$。
\end{solution}


\begin{example}
计算
\[
D_n=
\begin{vmatrix}
\sin(\alpha_1+\beta_1)&\sin(\alpha_1+\beta_2)&\cdots&\sin(\alpha_1+\beta_n)\\
\sin(\alpha_2+\beta_1)&\sin(\alpha_2+\beta_2)&\cdots&\sin(\alpha_2+\beta_n)\\
\vdots&\vdots&&\vdots\\
\sin(\alpha_n+\beta_1)&\sin(\alpha_n+\beta_2)&\cdots&\sin(\alpha_n+\beta_n)
\end{vmatrix}.
\]
\end{example}
\exampleblank

\begin{solution}
由
\[
\sin(\alpha_i+\beta_j)
=
\sin\alpha_i\cos\beta_j
+
\cos\alpha_i\sin\beta_j
\]
可知该矩阵是两个秩一矩阵之和，故秩不超过 $2$。因此当 $n\geq3$ 时
\[
\boxed{D_n=0}.
\]
当 $n=2$ 时，
\[
D_2
=
\sin(\alpha_1-\alpha_2)\sin(\beta_2-\beta_1).
\]
\end{solution}


\begin{example}
设 $A=(a_{ij})$，
\[
a_{ij}=\dfrac{1}{1-\alpha_i\beta_j},
\]
求 $|A|$。
\end{example}
\exampleblank

\begin{solution}
这是 Cauchy 型行列式。由标准 Cauchy 行列式公式，
\[
\boxed{
|A|
=
\dfrac{
\displaystyle
\prod\limits_{1\leq i<j\leq n}
(\alpha_i-\alpha_j)(\beta_i-\beta_j)
}{
\displaystyle
\prod\limits_{i,j=1}^{n}(1-\alpha_i\beta_j)
}
}.
\]
这里假设各分母非零；当某些 $\alpha_i$ 或 $\beta_j$ 重合时，公式自然给出行列式为零。
\end{solution}


\begin{example}
记
\[
s_k=\sum\limits_{j=1}^{n}a_j^k,
\qquad
s_0=n.
\]
计算
\[
D_n=
\begin{vmatrix}
s_0&s_1&s_2&\cdots&s_{n-1}\\
s_1&s_2&s_3&\cdots&s_n\\
s_2&s_3&s_4&\cdots&s_{n+1}\\
\vdots&\vdots&\vdots&&\vdots\\
s_{n-1}&s_n&s_{n+1}&\cdots&s_{2n-2}
\end{vmatrix}.
\]
\end{example}
\exampleblank

\begin{solution}
令 Vandermonde 矩阵
\[
V=
\begin{pmatrix}
1&1&\cdots&1\\
a_1&a_2&\cdots&a_n\\
\vdots&\vdots&&\vdots\\
a_1^{n-1}&a_2^{n-1}&\cdots&a_n^{n-1}
\end{pmatrix}.
\]
则题中矩阵正是
\[
VV^{\mathsf T}.
\]
因此
\[
D_n
=
(\det V)^2
=
\boxed{
\prod\limits_{1\leq i<j\leq n}(a_j-a_i)^2
}.
\]
\end{solution}


\begin{example}
计算
\[
D_n=
\begin{vmatrix}
1+x&x&\cdots&x\\
x&2+x&\cdots&x\\
\vdots&\vdots&\ddots&\vdots\\
x&x&\cdots&n+x
\end{vmatrix}.
\]
\end{example}
\exampleblank

\begin{solution}
仍由矩阵行列式引理，
\[
\boxed{
D_n
=
n!
\left(
1+x\sum\limits_{k=1}^{n}\dfrac{1}{k}
\right)
}.
\]
\end{solution}


\begin{example}
计算
\[
D_n=
\begin{vmatrix}
a_1+b_1&a_2&a_3&\cdots&a_n\\
a_1&a_2+b_2&a_3&\cdots&a_n\\
a_1&a_2&a_3+b_3&\cdots&a_n\\
\vdots&\vdots&\vdots&\ddots&\vdots\\
a_1&a_2&a_3&\cdots&a_n+b_n
\end{vmatrix}.
\]
\end{example}
\exampleblank

\begin{solution}
由秩一扰动公式，
\[
\boxed{
D_n
=
\prod\limits_{i=1}^{n}b_i
+
\sum\limits_{i=1}^{n}
a_i
\prod_{\substack{1\leq j\leq n\\j\neq i}}b_j
}.
\]
\end{solution}


\begin{example}
计算
\[
D_n=
\begin{vmatrix}
a&b&\cdots&b\\
b&a&\cdots&b\\
\vdots&\vdots&\ddots&\vdots\\
b&b&\cdots&a
\end{vmatrix}.
\]
\end{example}
\exampleblank

\begin{solution}
矩阵为
\[
(a-b)I_n+bJ,
\]
故
\[
\boxed{
D_n=(a-b)^{n-1}\bigl(a+(n-1)b\bigr)
}.
\]
\end{solution}


\begin{example}
计算
\[
D_n=
\begin{vmatrix}
a_1^2+1&a_1a_2&\cdots&a_1a_n\\
a_2a_1&a_2^2+1&\cdots&a_2a_n\\
\vdots&\vdots&\ddots&\vdots\\
a_na_1&a_na_2&\cdots&a_n^2+1
\end{vmatrix}.
\]
\end{example}
\exampleblank

\begin{solution}
题中矩阵为
\[
I_n+aa^{\mathsf T},
\qquad
a=(a_1,\ldots,a_n)^{\mathsf T}.
\]
故
\[
\boxed{
D_n=1+a^{\mathsf T}a
=
1+\sum\limits_{i=1}^{n}a_i^2
}.
\]
\end{solution}


\begin{example}
计算
\[
D_n=
\begin{vmatrix}
1+a_1+b_1&a_1+b_2&\cdots&a_1+b_n\\
a_2+b_1&1+a_2+b_2&\cdots&a_2+b_n\\
\vdots&\vdots&\ddots&\vdots\\
a_n+b_1&a_n+b_2&\cdots&1+a_n+b_n
\end{vmatrix}.
\]
\end{example}
\exampleblank

\begin{solution}
写成
\[
A=I_n+a\boldsymbol{1}^{\mathsf T}+\boldsymbol{1}b^{\mathsf T}.
\]
令
\[
U=(a,\boldsymbol{1}),
\qquad
V=(\boldsymbol{1},b).
\]
由
\[
\det(I_n+UV^{\mathsf T})
=
\det(I_2+V^{\mathsf T}U),
\]
得到
\[
\boxed{
D_n
=
\left(1+\sum\limits_{i=1}^{n}a_i\right)
\left(1+\sum\limits_{i=1}^{n}b_i\right)
-
n\sum\limits_{i=1}^{n}a_ib_i
}.
\]
\end{solution}


\begin{example}
计算
\[
D_n=
\begin{vmatrix}
0&a_1+a_2&\cdots&a_1+a_n\\
a_2+a_1&0&\cdots&a_2+a_n\\
\vdots&\vdots&\ddots&\vdots\\
a_n+a_1&a_n+a_2&\cdots&0
\end{vmatrix}.
\]
\end{example}
\exampleblank

\begin{solution}
先设各 $a_i\neq0$。记
\[
a=(a_1,\ldots,a_n)^{\mathsf T},
\qquad
D=-2\diag(a_1,\ldots,a_n).
\]
则
\[
A=D+a\boldsymbol{1}^{\mathsf T}+\boldsymbol{1}a^{\mathsf T}.
\]
对该秩二扰动使用矩阵行列式引理，可得
\[
\boxed{
D_n
=
(-1)^n2^{n-2}
\left(\prod\limits_{i=1}^{n}a_i\right)
\left[
(n-2)^2
-
\left(\sum\limits_{i=1}^{n}a_i\right)
\left(\sum\limits_{i=1}^{n}\dfrac{1}{a_i}\right)
\right]
}.
\]
两端均可视为 $a_i$ 的多项式恒等式，因此有 $a_i=0$ 的情形可由连续性或展开后代入得到。
\end{solution}


\begin{example}
计算
\[
D_n=
\begin{vmatrix}
\lambda&2&3&\cdots&n\\
1&\lambda+1&3&\cdots&n\\
1&2&\lambda+2&\cdots&n\\
\vdots&\vdots&\vdots&\ddots&\vdots\\
1&2&3&\cdots&\lambda+n-1
\end{vmatrix}.
\]
\end{example}
\exampleblank

\begin{solution}
记
\[
v=(1,2,\ldots,n)^{\mathsf T}.
\]
题中矩阵可写成
\[
A=(\lambda-1)I_n+\boldsymbol{1}v^{\mathsf T}.
\]
故
\[
\boxed{
D_n
=
(\lambda-1)^{n-1}
\left(
\lambda-1+\dfrac{n(n+1)}{2}
\right)
}.
\]
\end{solution}


\begin{example}
计算
\[
D_n=
\begin{vmatrix}
1+a_1&1&\cdots&1\\
1&1+a_2&\cdots&1\\
\vdots&\vdots&\ddots&\vdots\\
1&1&\cdots&1+a_n
\end{vmatrix}.
\]
\end{example}
\exampleblank

\begin{solution}
有
\[
A=\diag(a_1,\ldots,a_n)+J.
\]
故
\[
\boxed{
D_n
=
\prod\limits_{i=1}^{n}a_i
+
\sum\limits_{i=1}^{n}
\prod_{\substack{1\leq j\leq n\\j\neq i}}a_j
}.
\]
若各 $a_i\neq0$，也可写成
\[
D_n
=
\left(\prod\limits_{i=1}^{n}a_i\right)
\left(
1+\sum\limits_{i=1}^{n}\dfrac{1}{a_i}
\right).
\]
\end{solution}


\begin{example}[\markstar]
设 $A_{m\times n}$，$B_{n\times m}$，证明
\[
\lambda^n
\left|\lambda I_m-AB\right|
=
\lambda^m
\left|\lambda I_n-BA\right|.
\]
\end{example}
\exampleblank

\begin{solution}
先设 $\lambda\neq0$。由 Sylvester 行列式恒等式
\[
\det(I_m-XY)=\det(I_n-YX),
\]
取
\[
X=\dfrac{1}{\lambda}A,\qquad Y=B,
\]
得
\[
\det\left(I_m-\dfrac{1}{\lambda}AB\right)
=
\det\left(I_n-\dfrac{1}{\lambda}BA\right).
\]
乘以 $\lambda^{m+n}$ 即得
\[
\lambda^n
\det(\lambda I_m-AB)
=
\lambda^m
\det(\lambda I_n-BA).
\]
两边均为 $\lambda$ 的多项式，故恒等式对 $\lambda=0$ 也成立。
\end{solution}


\newpage

\section{综合问题}

\begin{example}
证明下面的行列式不为零：
\[
D_n=
\det\left((i+j-1)^i\right)_{1\leq i,j\leq n}.
\]
原题取 $n=2011$。
\end{example}
\exampleblank

\begin{solution}
令
\[
x_j=j-1,\qquad
p_i(x)=(x+i)^i.
\]
则
\[
D_n=\det\bigl(p_i(x_j)\bigr)_{1\leq i,j\leq n}.
\]
将
\[
p_i(x)
=
\sum\limits_{k=0}^{i}
C_k^i i^{\,i-k}x^k
\]
看作单项式基下的系数展开。其常数项与一次项系数均为 $i^i$，而最高次项系数为 $1$。用 Cauchy--Binet 公式展开后，只有“删去 $x^0$ 列”与“删去 $x^1$ 列”两个系数子式可能非零，且二者均为 $1$。

于是
\[
D_n
=
\Delta(x_1,\ldots,x_n)
\bigl(e_n(x_1,\ldots,x_n)+e_{n-1}(x_1,\ldots,x_n)\bigr).
\]
这里
\[
x_1=0,x_2=1,\ldots,x_n=n-1,
\]
故
\[
e_n=0,\qquad
e_{n-1}=(n-1)!,
\]
并且
\[
\Delta(0,1,\ldots,n-1)
=
\prod\limits_{k=1}^{n-1}k!.
\]
因此
\[
\boxed{
D_n
=
(n-1)!
\prod\limits_{k=1}^{n-1}k!
>0
}.
\]
特别地，原题 $n=2011$ 时行列式必不为零。
\end{solution}


\begin{example}[\markstar]
计算 Cauchy 行列式
\[
D_n=
\begin{vmatrix}
\dfrac{1}{a_1+b_1}&\dfrac{1}{a_1+b_2}&\cdots&\dfrac{1}{a_1+b_n}\\[6pt]
\dfrac{1}{a_2+b_1}&\dfrac{1}{a_2+b_2}&\cdots&\dfrac{1}{a_2+b_n}\\
\vdots&\vdots&&\vdots\\[4pt]
\dfrac{1}{a_n+b_1}&\dfrac{1}{a_n+b_2}&\cdots&\dfrac{1}{a_n+b_n}
\end{vmatrix}.
\]
\end{example}
\exampleblank

\begin{solution}
假设各 $a_i+b_j\neq0$。Cauchy 行列式公式为
\[
\boxed{
D_n
=
\dfrac{
\displaystyle
\prod\limits_{1\leq i<j\leq n}(a_j-a_i)(b_j-b_i)
}{
\displaystyle
\prod\limits_{i,j=1}^{n}(a_i+b_j)
}
}.
\]
可由对 $n$ 归纳，或逐步作列差后提取因子证明。
\end{solution}


\begin{example}[\markstar]
设 $n$ 阶行列式 $\det A$ 的元素都是关于 $t$ 的可微函数，证明：
\begin{enumerate}
  \item
  \[
  \dfrac{\mathrm{d}}{\mathrm{d}t}\det A
  =
  \det A_1+\cdots+\det A_n,
  \]
  其中 $A_i$ 为仅对 $A$ 的第 $i$ 行逐项求导所得矩阵；
  \item
  \[
  \dfrac{\mathrm{d}}{\mathrm{d}t}\det A
  =
  \sum\limits_{i,j=1}^{n}
  \dfrac{\mathrm{d}a_{ij}(t)}{\mathrm{d}t}A_{ij}.
  \]
\end{enumerate}
\end{example}
\exampleblank

\begin{solution}
行列式关于各行分别线性。对
\[
\det(R_1(t),\ldots,R_n(t))
\]
求导，按乘积法则的多线性版本得到
\[
\dfrac{\mathrm{d}}{\mathrm{d}t}\det A
=
\sum\limits_{i=1}^{n}
\det(R_1,\ldots,R_i',\ldots,R_n),
\]
即第一式。

再对每个 $\det A_i$ 按第 $i$ 行展开，
\[
\det A_i
=
\sum\limits_{j=1}^{n}a_{ij}'(t)A_{ij}.
\]
对 $i$ 求和即得第二式。
\end{solution}


\begin{example}[\markstar]
设 $A=(a_{ij})\neq0$，且对一切 $i,j$ 有
\[
a_{ij}=A_{ij},
\]
其中 $A_{ij}$ 为 $a_{ij}$ 的代数余子式。求 $|A|$。
\end{example}
\exampleblank

\begin{solution}
设 $A$ 为实 $n$ 阶矩阵，并先取 $n\geq3$。由于代数余子式矩阵等于 $A$，故
\[
\operatorname{adj}(A)=A^{\mathsf T}.
\]
于是
\[
AA^{\mathsf T}
=
A\operatorname{adj}(A)
=
|A|I_n.
\]
因 $A\neq0$ 且左端半正定，若 $|A|=0$ 则
\[
AA^{\mathsf T}=0
\]
将推出 $A=0$，矛盾。因此 $|A|>0$。取行列式得
\[
|A|^2=|A|^n,
\]
故
\[
|A|^{n-2}=1.
\]
当 $n\geq3$ 时
\[
\boxed{|A|=1}.
\]

原题若允许 $n=2$，则结论并不唯一：此时
\[
A=
\begin{pmatrix}
a&b\\
-b&a
\end{pmatrix},
\qquad
|A|=a^2+b^2,
\]
可取任意正值。因此原题需要补充 $n\geq3$ 才能得到唯一答案。
\end{solution}


\begin{example}
设 $a_1,a_2,\ldots,a_n$ 是 $n$ 个 $n$ 位自然数，其中 $a_i$ 的个位、十位、$\ldots$、第 $n$ 位数字依次为
\[
a_{i1},a_{i2},\ldots,a_{in},
\]
并设
\[
d=
\begin{vmatrix}
a_{11}&a_{12}&\cdots&a_{1n}\\
a_{21}&a_{22}&\cdots&a_{2n}\\
\vdots&\vdots&&\vdots\\
a_{n1}&a_{n2}&\cdots&a_{nn}
\end{vmatrix}.
\]
证明：$a_1,a_2,\ldots,a_n$ 的最大公因数 $r$ 整除 $d$。
\end{example}
\exampleblank

\begin{solution}
由十进制展开，
\[
a_i
=
a_{i1}+10a_{i2}+\cdots+10^{n-1}a_{in}.
\]
对行列式作列变换
\[
C_1
\longleftarrow
C_1+10C_2+\cdots+10^{n-1}C_n.
\]
行列式的值不变，而新第一列变为
\[
(a_1,a_2,\ldots,a_n)^{\mathsf T}.
\]
因为 $r$ 整除每个 $a_i$，故可从第一列提出因子 $r$，于是
\[
\boxed{r\mid d}.
\]
\end{solution}


\begin{example}
设
\[
D_n=
\begin{vmatrix}
1&-1&-1&\cdots&-1\\
1&1&-1&\cdots&-1\\
1&1&1&\cdots&-1\\
\vdots&\vdots&\vdots&\ddots&\vdots\\
1&1&1&\cdots&1
\end{vmatrix}.
\]
求 $D_n$ 展开式中正项的项数。
\end{example}
\exampleblank

\begin{solution}
由于矩阵中没有零元，Leibniz 展开共有 $n!$ 项。设正项数为 $P$，负项数为 $N$，则
\[
P+N=n!,
\qquad
P-N=D_n.
\]
从最后一行开始依次作
\[
R_i\longleftarrow R_i-R_{i-1},
\qquad i=n,n-1,\ldots,2,
\]
得到除第一行外每行仅在第 $i$ 列有元素 $2$，故
\[
D_n=2^{n-1}.
\]
因此
\[
\boxed{
P=\dfrac{n!+2^{n-1}}{2}
}.
\]
\end{solution}


\begin{example}
计算
\[
\sum_{j_1j_2\cdots j_n}
\begin{vmatrix}
a_{1j_1}&a_{1j_2}&\cdots&a_{1j_n}\\
a_{2j_1}&a_{2j_2}&\cdots&a_{2j_n}\\
\vdots&\vdots&&\vdots\\
a_{nj_1}&a_{nj_2}&\cdots&a_{nj_n}
\end{vmatrix},
\]
其中求和遍历 $1,2,\ldots,n$ 的所有全排列
\[
j_1j_2\cdots j_n.
\]
\end{example}
\exampleblank

\begin{solution}
每个行列式都是原矩阵 $A=(a_{ij})$ 的列置换，因此对置换 $\sigma$，
\[
\det A_\sigma
=
\operatorname{sgn}(\sigma)\det A.
\]
当 $n\geq2$ 时，偶置换与奇置换数目相同，故
\[
\boxed{
\sum_{\sigma\in S_n}\det A_\sigma=0
}.
\]
$n=1$ 时结果为 $a_{11}$。
\end{solution}


\begin{example}[\markstar]
设 $n\geq3$，
\[
f_1(x),f_2(x),\ldots,f_n(x)
\]
都是次数不超过 $n-2$ 的多项式，$a_1,a_2,\ldots,a_n$ 为任意数。证明
\[
\begin{vmatrix}
f_1(a_1)&f_2(a_1)&\cdots&f_n(a_1)\\
f_1(a_2)&f_2(a_2)&\cdots&f_n(a_2)\\
\vdots&\vdots&&\vdots\\
f_1(a_n)&f_2(a_n)&\cdots&f_n(a_n)
\end{vmatrix}
=0.
\]
\end{example}
\exampleblank

\begin{solution}
每个 $f_j$ 都属于
\[
\Span\{1,x,\ldots,x^{n-2}\},
\]
该多项式空间维数为 $n-1$。于是 $n$ 个多项式
\[
f_1,\ldots,f_n
\]
线性相关。即存在不全为零的 $\lambda_j$ 使
\[
\sum\limits_{j=1}^{n}\lambda_jf_j(x)\equiv0.
\]
代入 $x=a_i$，得到题中矩阵的列线性相关，故行列式为零。
\end{solution}


\begin{example}
计算
\[
D=
\begin{vmatrix}
1^{50}&2^{50}&3^{50}&\cdots&100^{50}\\
2^{50}&3^{50}&4^{50}&\cdots&101^{50}\\
3^{50}&4^{50}&5^{50}&\cdots&102^{50}\\
\vdots&\vdots&\vdots&&\vdots\\
100^{50}&101^{50}&102^{50}&\cdots&199^{50}
\end{vmatrix}.
\]
\end{example}
\exampleblank

\begin{solution}
第 $j$ 列的第 $i$ 个元素为
\[
(i+j-1)^{50}.
\]
把它看成关于 $i$ 的多项式，则所有列都属于由
\[
1,i,i^2,\ldots,i^{50}
\]
张成的 $51$ 维空间。因此 $100$ 个列向量必线性相关，故
\[
\boxed{D=0}.
\]
\end{solution}


\begin{proposition}[\marktwostar\ 循环矩阵]
设
\[
I=
\begin{pmatrix}
0&1&0&\cdots&0\\
0&0&1&\cdots&0\\
\vdots&\vdots&\vdots&&\vdots\\
0&0&0&\cdots&1\\
1&0&0&\cdots&0
\end{pmatrix},
\qquad
J=
\begin{pmatrix}
0&1&0&\cdots&0\\
0&0&1&\cdots&0\\
\vdots&\vdots&\vdots&&\vdots\\
0&0&0&\cdots&1\\
0&0&0&\cdots&0
\end{pmatrix}.
\]
设
\[
A=
\begin{pmatrix}
a_0&a_1&a_2&\cdots&a_{n-1}\\
a_{n-1}&a_0&a_1&\cdots&a_{n-2}\\
\vdots&\vdots&\vdots&&\vdots\\
a_2&a_3&a_4&\cdots&a_1\\
a_1&a_2&a_3&\cdots&a_0
\end{pmatrix},
\]
并令
\[
f(x)=a_0+a_1x+\cdots+a_{n-1}x^{n-1}.
\]
若 $\omega_1,\ldots,\omega_n$ 为全部 $n$ 次单位根，则：
\begin{enumerate}
  \item $I^n=E$，$J^n=0$，且 $\omega_i$ 为 $I$ 的特征值；
  \item 矩阵 $B$ 与 $I$ 可交换当且仅当 $B=g(I)$；矩阵 $B$ 与 $J$ 可交换当且仅当 $B=g(J)$，其中 $g$ 为次数小于 $n$ 的多项式；
  \item $f(\omega_i)$ 为 $A$ 的全部特征值，且
  \[
  |A|=\prod\limits_{i=1}^{n}f(\omega_i).
  \]
\end{enumerate}
\end{proposition}

\begin{proof}
前两式 $I^n=E,J^n=0$ 由移位矩阵的定义直接得到。对任意 $n$ 次单位根 $\omega$，
\[
v(\omega)
=
(1,\omega,\omega^2,\ldots,\omega^{n-1})^{\mathsf T}
\]
满足
\[
Iv(\omega)=\omega v(\omega),
\]
故全部 $n$ 次单位根恰为 $I$ 的特征值。

若 $BI=IB$，比较矩阵元素可知 $B$ 的各行由上一行循环移位得到，因此 $B$ 为循环矩阵，即
\[
B=b_0E+b_1I+\cdots+b_{n-1}I^{n-1}=g(I).
\]
反之任意 $g(I)$ 显然与 $I$ 可交换。

类似地，$BJ=JB$ 强迫 $B$ 为上三角 Toeplitz 矩阵，从而
\[
B=b_0E+b_1J+\cdots+b_{n-1}J^{n-1}=g(J).
\]

最后
\[
A=f(I),
\]
故若 $Iv(\omega_i)=\omega_i v(\omega_i)$，则
\[
Av(\omega_i)
=
f(I)v(\omega_i)
=
f(\omega_i)v(\omega_i).
\]
因此
\[
|A|
=
\prod\limits_{i=1}^{n}f(\omega_i).
\]
\end{proof}


\begin{example}
计算：
\begin{enumerate}
  \item
  \[
  \begin{vmatrix}
  5&6&7&8\\
  6&7&8&5\\
  7&8&5&6\\
  8&5&6&7
  \end{vmatrix};
  \]
  \item
  \[
  D_n=
  \begin{vmatrix}
  1&2&\cdots&n-1&n\\
  2&3&\cdots&n&1\\
  \vdots&\vdots&&\vdots&\vdots\\
  n&1&\cdots&n-2&n-1
  \end{vmatrix}.
  \]
\end{enumerate}
\end{example}
\exampleblank

\begin{solution}
第一问由例题 2.2.3 已知
\[
\boxed{-416}.
\]

第二问是左循环行列式。先记 $C_n$ 为第一行同为
\[
(1,2,\ldots,n)
\]
的标准右循环行列式。令
\[
f(x)=1+2x+\cdots+nx^{n-1}.
\]
当 $\omega^n=1,\omega\neq1$ 时，
\[
f(\omega)
=
\sum\limits_{k=1}^{n}k\omega^{k-1}
=
\dfrac{n}{\omega-1},
\]
而
\[
f(1)=\dfrac{n(n+1)}{2}.
\]
又
\[
\prod_{\substack{\omega^n=1\\\omega\neq1}}(1-\omega)=n,
\]
故
\[
\det C_n
=
(-1)^{n-1}
\dfrac{(n+1)n^{n-1}}{2}.
\]

题中左循环矩阵可由 $C_n$ 作行置换
\[
0,1,\ldots,n-1
\longmapsto
0,n-1,n-2,\ldots,1
\]
得到，该置换由
\[
\left\lfloor\dfrac{n-1}{2}\right\rfloor
\]
个对换组成。因此
\[
\boxed{
D_n
=
(-1)^{\,n-1+\left\lfloor\frac{n-1}{2}\right\rfloor}
\dfrac{(n+1)n^{n-1}}{2}
}.
\]
\end{solution}


\begin{example}
计算
\[
D_n=
\begin{vmatrix}
a_0&a_1&a_2&\cdots&a_{n-1}\\
\mu a_{n-1}&a_0&a_1&\cdots&a_{n-2}\\
\mu a_{n-2}&\mu a_{n-1}&a_0&\cdots&a_{n-3}\\
\vdots&\vdots&\vdots&\ddots&\vdots\\
\mu a_1&\mu a_2&\mu a_3&\cdots&a_0
\end{vmatrix}.
\]
\end{example}
\exampleblank

\begin{solution}
设
\[
T=
\begin{pmatrix}
0&1&0&\cdots&0\\
0&0&1&\cdots&0\\
\vdots&\vdots&\vdots&&\vdots\\
0&0&0&\cdots&1\\
\mu&0&0&\cdots&0
\end{pmatrix},
\]
则
\[
T^n=\mu I_n,
\qquad
A=a_0I_n+a_1T+\cdots+a_{n-1}T^{n-1}.
\]
取 $\rho$ 满足 $\rho^n=\mu$，并令 $\omega_k$ 遍历全部 $n$ 次单位根，则 $T$ 的特征值为
\[
\rho\omega_k.
\]
因此
\[
\boxed{
D_n
=
\prod\limits_{k=1}^{n}
\left(
a_0+a_1\rho\omega_k+\cdots+a_{n-1}(\rho\omega_k)^{n-1}
\right)
}.
\]
等价地，乘积遍历所有满足 $z^n=\mu$ 的 $z$：
\[
D_n
=
\prod_{z^n=\mu}
\left(
a_0+a_1z+\cdots+a_{n-1}z^{n-1}
\right).
\]
\end{solution}


\begin{example}
计算
\[
D_n=
\begin{vmatrix}
a_1&a_2&a_3&\cdots&a_n\\
-a_n&a_1&a_2&\cdots&a_{n-1}\\
-a_{n-1}&-a_n&a_1&\cdots&a_{n-2}\\
\vdots&\vdots&\vdots&\ddots&\vdots\\
-a_2&-a_3&-a_4&\cdots&a_1
\end{vmatrix}.
\]
\end{example}
\exampleblank

\begin{solution}
这是上一题取 $\mu=-1$ 的情形。令
\[
f(z)=a_1+a_2z+\cdots+a_nz^{n-1}.
\]
则
\[
\boxed{
D_n
=
\prod_{z^n=-1}f(z)
=
\prod_{z^n=-1}
\left(
a_1+a_2z+\cdots+a_nz^{n-1}
\right)
}.
\]
\end{solution}
